\scanpage{41}\mspage{40}
% L1–3
\uline{Caractérisation d'un} \uline{intervalle} \deleted{réseau-interval} \uline{ordonné}, \uline{en \deleted{tant} que réseau}

% L4–5
\uline{par la donnée} de
\diagram{\begin{tikzcd}[ampersand replacement=\&,row sep=0.75em,column sep=0.15em]
S \& {}\subset\unsure{\mathcal P|\mathscr P|\mathrm P}_2(L)\\
\makebox[0pt][c]{\(\operatorname{Drap}_2(L)\)}\arrow[u,equal] \&
\end{tikzcd}}
% L6–9
et pour $\varepsilon=\{a,b\}\in S$, la décomposition de \deleted{$L$ en} $\complement_{\mathrm{mod}}\,\varepsilon=\{x\unsure{+|\in}L\mid\{\unsure{a;x|a',x|a,x}\},\{b,x\}\in S\}$ en ``composantes connexes'' (par relation $R_\varepsilon$)

% L10–18
J'ai commencé à faire l'exercice, mais j'ai l'impression que ça n'en vaut pas la peine de l'expliciter jusqu'au bout. Ça doit avoir été fait depuis longtemps (Hilbert~?) dans le cas totalement ordonné \inserted{(et divisible)} (en termes \unsure{tout au|tant au|tout aux} moins de la relation~: $s$ est entre $a$ et $b$ \ldots\ )
